\section{Related work}

Skip this when you read this script for the first time


There are two very important pieces of related work:
The first one is tools like Intel's Application Performance Snapshot\footnote{See for example \url{https://www.intel.com/content/www/us/en/docs/vtune-profiler/user-guide-application-snapshot-linux}} (APS) or the ClusterCockpit\footnote{\url{https://clustercockpit.org}} that \emph{monitor} the code's behaviour and deliver a high-level overview how a machine uses the underlying system.
Rather than piping a code through a set of assessment steps, these tools give us a live fingerprint of a code's behaviour.
These tools still fail to uncover correlations between the different performance dimensions, but we might argue that they yield more realistic insights under ``working conditions'', whereas we follow more of a lab-based procedure.


The second idea to dig into the root of all evil by successively focusing on refined metrics, ruling out factors that are not relevant for the quantities of interest, is not ours.
The Performance Optimisation and Productivity (POP) consortium\footnote{\url{https://pop-coe.eu}}---a EU Centre of Excellence in HPC---has published some very detailed guidelines for this already, and the members of this consortium have also worked on tools that help us to assess code.
Several ideas in the present chapter are heavily inspired by these POP metrics.
The collection here is my braindump.
It certainly is highly recommended to study the material and metrics of POP.
I'll indeed frequently refer to their procedures, tools and recommendations.

